% Appendix with required statements and checklists for Agents4Science 2025

% Provide lightweight fallbacks for checklist macros if the style doesn't define them
\providecommand{\answerYes}{Yes}
\providecommand{\answerNo}{No}
\providecommand{\answerNA}{N/A}

\providecommand{\involvementA}{Human-generated}
\providecommand{\involvementB}{Mostly human, assisted by AI}
\providecommand{\involvementC}{Mostly AI, assisted by human}
\providecommand{\involvementD}{AI-generated}

\section*{Responsible AI Statement}
This research adheres to community ethical guidelines and promotes responsible AI development practices. The work focuses on cryptographic algorithm analysis without creating or enabling malicious applications.

\textbf{Positive societal impacts:}
\begin{itemize}
  \item Provides empirical evidence for secure hash function selection in security-critical applications
  \item Contributes to the understanding of cryptographic algorithm performance characteristics
  \item Enables informed decision-making for cryptographic system design
  \item Supports the development of more secure and efficient cryptographic systems
\end{itemize}

\textbf{Potential negative impacts and mitigations:}
\begin{itemize}
  \item Analysis results could be misused to identify vulnerabilities; we emphasize adoption of modern, secure algorithms and explicitly discourage use of deprecated ones (e.g., MD5)
  \item Benchmarks could be abused to optimize attacks on weaker algorithms; we provide context and caveats and do not publish novel attack vectors
  \item All evaluated algorithms and settings are public and well-studied; no sensitive data or systems were targeted
\end{itemize}

\section*{Reproducibility Statement}
We release code, configuration, and generated artifacts to facilitate reproduction. Experiments were run on a CPU-only macOS system (Darwin 24.6.0) with Python 3.13. The software stack is pinned in \texttt{code/requirements.txt} (e.g., numpy 2.3.3, matplotlib 3.10.6, seaborn 0.13.2). The repository includes:\
\noindent\textbf{Code}: \texttt{code/} (analysis and runner).\
\noindent\textbf{Data}: synthetic test generators in code, with outputs written to \texttt{results/}.\
\noindent\textbf{Artifacts}: summary files in \texttt{results/}, and figures in \texttt{results/figures/}.\
We specify input sizes, input types (random/structured/edge-case), number of trials per configuration, and timing methodology in the Methodology section. To reproduce: create a Python 3.13 environment, install \texttt{requirements.txt}, and run the provided runner script. We report means and include dispersion recommendations; confidence intervals can be produced by re-running with multiple seeds.

\textbf{Reproducibility note}: Results reported in this paper were generated with random seed 2025 (set via \texttt{SEED=2025} environment variable). The analysis script accepts a \texttt{SEED} environment variable to ensure deterministic generation of test vectors. To reproduce exact results, run with \texttt{SEED=2025}; for different random samples, use a different seed value.

\section*{Agents4Science AI Contribution Disclosure}
\textbf{Hypothesis development}: \involvementB{}\\
Briefly: Humans defined the scope (hash functions, metrics); AI assisted with drafting and editing.

\textbf{Experimental design and implementation}: \involvementB{}\\
Briefly: Humans implemented and reviewed the code; AI provided refactoring suggestions and documentation edits.

\textbf{Analysis of data and interpretation of results}: \involvementB{}\\
Briefly: Humans conducted the analysis and validated findings; AI assisted with figure caption wording and summarization.

\textbf{Writing}: \involvementB{}\\
Briefly: The manuscript was primarily written by humans with AI assistance for wording, organization, and grammar.

\section*{Agents4Science Paper Checklist}
% Answers use simple macros; replace with more granular discussion as needed
\begin{enumerate}
  \item \textbf{Claims}: Do the abstract and introduction accurately reflect contributions and scope?\\
  Answer: \answerYes{} — Claims match methods and reported results; limitations are discussed.

  \item \textbf{Limitations}: Are limitations discussed?\\
  Answer: \answerYes{} — See the Limitations subsection in Discussion.

  \item \textbf{Theory assumptions and proofs}: Are assumptions/proofs complete (if applicable)?\\
  Answer: \answerNA{} — This work is empirical; no new formal theorems are introduced.

  \item \textbf{Experimental reproducibility}: Is sufficient information provided to reproduce key results?\\
  Answer: \answerYes{} — Code, parameters, data generation process, and environment details are provided.

  \item \textbf{Open access to data and code}: Is code/data available with instructions?\\
  Answer: \answerYes{} — Scripts and instructions are included in \texttt{code/} with version-pinned dependencies.

  \item \textbf{Experimental details}: Are train/test details and settings specified?\\
  Answer: \answerYes{} — Input sizes, data types, iterations per configuration, and timing methods are specified.

  \item \textbf{Statistical significance}: Are dispersion or uncertainty measures reported or supported?\\
  Answer: \answerYes{} — Summary statistics are provided; the code supports computing confidence intervals via repeated runs.

  \item \textbf{Compute resources}: Are compute resources and runtime described?\\
  Answer: \answerYes{} — CPU-only macOS system; runtimes and vector counts are reported in Methodology.

  \item \textbf{Code of ethics}: Does the work conform to the conference code of ethics?\\
  Answer: \answerYes{} — The study evaluates public algorithms and emphasizes secure usage recommendations.

  \item \textbf{Broader impacts}: Are potential societal impacts discussed?\\
  Answer: \answerYes{} — See Responsible AI Statement.
\end{enumerate}




